\chapter{Prerequisiti}\stepcounter{capitoloeccolo}

\section{Prerequisiti di Logica}

\lipsum[1]

\subsection{Assioma della scelta}

Per la costruzione di $ \starR $ come ultraprodotto, sarà necessario utilizzare l'Assioma della Scelta. Lo si presenta anche nella sua forma equivalente: il Lemma di Zorn

\paragrafo{Assioma di scelta}{%
    Per ogni insieme $ X\neq \emptyset $ esiste \[f:\parti{X}\setminus\{\emptyset\}\longrightarrow X\]tale che $ f(Y) \in Y $ per ogni $ Y \subseteq X $ non vuoto.
}{}{}
\paragrafo{Lemma di Zorn}{%
    Se un insieme (parzialmente) ordinato è tale che ogni catena ha un maggiorante, allora ha un elemento massimale.
}{}{}
\proposizione{odijcoijodijocijoijoijoi}{
    Sia $ F:\mathcal{N}\longrightarrow \mathcal{M} $ un'immersione. Le seguenti affermazioni sono equivalenti: \begin{enumerate}
        \item $ F $ è un'immersione elementare; 
        \item per ogni $ \mathcal{L} $-formula $ \varphi(y,x_1,\dots,x_{n} ) $ ed ogni $ a_1,\dots,a_{n} \in N  $ \[
            \mathcal{M}\vDash \exists y \, \varphi\left[y, F(a_1),\dots,F(a_{n}) \right]\qquad\text{sse}\qquad \parbox{15em}{$\mathcal{M}\vDash\varphi\left[F(a_{0}), F(a_1),\dots,F(a_{n}) \right]${ per qualche }$a_0 \in N.$}
        \]
    \end{enumerate}
}

\section{Prerequisiti di Geometria Differenziale}\label{sec:geodiff}

\lipsum[2]